\section{Numerical choices}\label{sec:constants}

Fix the instruction set, finite track order, tape template and polynomial time allowance first. Choose \(w\) sufficiently large for every polynomial comparison with \(P_0\), every constant-times-\(A\) comparison with \(v\), and the collar inequalities. A lower bound of \(10^{10}\) is compatible with all specified finite lower bounds, and \(w\) may be increased to satisfy the interpreter allowance and the fixed collar requirements. All native lengths are then fixed. In particular \(d_1\ll P_0\), and any fixed polynomial in \(w\) is \(o(P_0)\). The estimates use fixed proof constants independent of word width. The following table lists the choices and the larger allowances used when estimates overlap.

\begin{longtable}{@{}>{\raggedright\arraybackslash}p{.30\linewidth}>{\raggedright\arraybackslash}p{.62\linewidth}@{}}
\toprule
Quantity & Value and role \\
\midrule
\endhead
\(q,A,B,v\) & \(w^2,w^3,128q,w^4\), respectively. They determine the propagation length, Shrink thresholds, provisional pass and Grow rounds. \\
\(P_0,P,T\) & \(2^w,5P_0\), and the duration of the fixed clock in Section~\ref{sec:clock}. \\
\(H_0,d_0,d_1\) & \(2q+2,(2A+4)H_0,qd_0\). These are the handle, Main cycle and active Shrink spacing. \\
\(r_{\mathrm{last}},s_0,s_1,s_c\) & Start of final withdrawal, start of ordinary phase, end of \(\operatorname{Shrink}(0)\), and the service date in early computation. \\
\(m,d\), incoming pad & \(m=\lfloor w/5\rfloor\), \(d=4w+8\), and \(64A\) child units. Common span \(3w+2m+4+2(608)<d\): the \(2\cdot608\) term is fixed slack; owner-level displacements are paid by the m-owner collar (\(608<m\)) and do not enlarge K. \\
\(M\), verifier threshold & \(M=128d\), count strictly greater than \(4d\), epoch length \(M+1\). \\
Vacant-band thresholds & \(4B+4\) at all retired strobes, \(8B+8\) before a partial pass. At most \(B\) positions can be entered from each side of a pass, so the central band stays vacant. Width \(4B+4\) works throughout the ordinary and terminal estimates. \\
End gap and old material & Use \(16B+2d+24\) in the early certificate and comparison with \(\Delta\). After removal of \(K\) and boundary positions it leaves a vacant band of width \(4B+4\); also \(\Delta<16B+2d+24<3\Delta\). \\
\(\Delta\) and the survivor budgets & \(\Delta=1024q\); the budgets are \(3A+3\Delta\) and \(6A+3\Delta\). The assignment allowance is \(9d+80\); surviving faces advance at least \(B-20\) in the partial pass. \\
Common full seed & Derived \(100A d_1\) bound per home and \(N=200A d_1\) as uniform search allowance. For \(b_{\rm seed}=6N+4w^6+20\), search \(b_{\rm seed}+1\) trials of \(5d_0+1\) dates. The transport margin adds the relaxation and preparation delays to this bound. \\
Source density & \(L=60B+3d+40\) for the vacant interval comparison, at most five labels and \(4d+12\) components per label. The resulting deficit is at most \(2{,}000{,}000A\). \\
Terminal bands & Thin radius \(100B\), near-source radius \(300B\), wide half-width \(v/5\). Permanent gain at most \(35B\) across the event. Delete at most \((220A+2w+20)d_0\) terminal comparison dates. \\
Endpoint repair & Search length \(3v d_0+2w^6+12d_0\), clear zone \(100q\), retained guard at depth below \(97q\). Use \(4q+4\) inspections for one search and \(8q+12\) for two. Strip length \(17q+25\) leaves a \(q+1\) backing after two-sided losses, with the nearest witness within \(47q\). Final common exception has width at most \(d+q+42<2q\). \\
Directed advance & Startup in \(6d_0\) clean ticks, allow \(12d_0\) with the pulse. Cost \(d+4(q+20)<5q\) old sites and six cycle equivalents; overshoot \(P+10v\). \\
Nested tests & \(M_0=\lfloor2w/5\rfloor\); fixed multipliers with four rounding units. The cap test pays the extra spatial factor \(4T/P+4\) and temporal factor \(6\), before its one-rank collar. \\
\(D_0,H_n\) & Fix \(D_0\) after \(P,T,w\) and test multipliers; set \(H_n=D_0^{n+2}\). Fix the paint multiplier at this stage. \\
\(M_s,b_h,v_h,s_h\) & Choose the required short exponent and then its refresh, replica and reader counts. Require \(b_h>20\) and strict exponent slack in Lemma~\ref{lem:stripes}. \\
\(M_l,m_h\) & After \(b_h\) fix the long admission exponent and choose \(m_h\) for the long estimate. Use odd reader and channel counts. \\
\(D_{\rm lane},\ell_0,n_0\) & Choose the stride and base word width after the helper counts. Increase \(n_0\) last for the uniform probability and time bounds on all \(I_n\). \\
\bottomrule
\end{longtable}


\subsection{Notation and local scopes}\label{sec:notation}
The following conventions distinguish global parameters from the explicitly local variables reused in individual calculations.
\begin{longtable}{@{}>{\raggedright\arraybackslash}p{.28\linewidth}>{\raggedright\arraybackslash}p{.64\linewidth}@{}}
\toprule
Notation & Meaning and scope \\
\midrule
\endhead
\(\Delta,\Delta_{\rm leaf}\) & \(\Delta=1024q\) is the ordinary depth margin; \(\Delta_{\rm leaf}\) is the enlarged-leaf overlap degree in Section~\ref{sec:joint}. \\
\(H,H_0,H_n\) & H denotes a spatial home in organization proofs, the suffix length \(\lceil p^{-6}\rceil\) throughout Section~\ref{sec:receiver}, and a cone height where explicitly stated. \(H_0=2q+2\) is always the handle; \(H_n=D_0^{n+2}\) is a stripe-box radius. The decoding exception is \(I_{\rm dec}\). \\
\(L,P_0\) & L is a local gap bound, strip length or time horizon, defined afresh in its lemma, the segment length in Sections~\ref{sec:receiver}--\ref{sec:stripes}, and the leaf count \(L=2^{i-j}\) in Section~\ref{sec:joint}; the Grow label L in the advance proof denotes the left donor. \(P_0=P/5\) is the packet length, including in the classification calculation. \\
\(D,D_{\rm tr}\) & D denotes the dead symbol in state expressions. In numerical proofs D is explicitly a pair count, a \(35B\) gain, a rear-wall interval, or the advance debit \(d+4(q+20)\). \(D_{\rm tr}\) is the ballot transport margin. \(D_0,D_j,D_{\rm lane}\) have their displayed geometric, recursive and stride meanings. \\
\(b_0,b_j,b_{\rm dst}\) & \(b_0=0,b_j\) are service lookahead. The destination coordinate in the ballot path equation is \(b_{\rm dst}\). The probability \(b_j\) in the iid sparsity proof is local to that proof. Other local b's denote a scale base, lock rank, interval end or epoch offset, as explicitly introduced. \\
\(W,W_0\) & W is \(4B+4\) in sterile-band proofs, \(\lceil p^{-3/2}\rceil\) in Sections~\ref{sec:receiver}--\ref{sec:stripes}, and the explicitly displayed end-search duration in the end ledger. \(W_0\) is the protected two-period test. \\
\(N,p,Y\) & N is a nonfull-date allowance, a local Grow-cycle count, or the alphabet size in Section~\ref{sec:gibbs}. The probe rate is p; a local endgap also called p is confined to ordinary geometry. The PCA probability table is \(p(a\mid\xi)\). Y is the native carrier; \(Y_i\) are autonomous reference channels. \\
\(v,w,\beta_-,\beta_+\) & \(v=w^4\) and w are native geometry parameters; v also denotes a space--time site in \eqref{eq:wsm} and Section~\ref{sec:gibbs}; \(\beta_-<\beta_+\) are rational inverse-temperature endpoints. \\
Nonfull & At least one of a home's P positions lacks an actual matching served child at that fine date. \\
\bottomrule
\end{longtable}

\section{Case tables for the numerical bounds}\label{sec:case-tables}

The rows below collect the timing reserves used in Lemmas~\ref{lem:advance}, \ref{lem:reader} and~\ref{lem:receiver-mixing}. The corresponding calculations are in the proofs of those lemmas. Each row uses the same native clock.

\subsection{Directed interruption}

\begin{longtable}{@{}>{\raggedright\arraybackslash}p{.25\linewidth}>{\raggedright\arraybackslash}p{.38\linewidth}>{\raggedright\arraybackslash}p{.28\linewidth}@{}}
\toprule
State at interruption & Repair & Timing and vacancy bounds \\
\midrule
\endhead
Startup & Restart from the \(K\)-avoiding ordinary backing and use a complete clear pass. & At most \(12d_0\) from the cut. Discard one cycle. \\
Free attack or free withdrawal & Compare to the last unmarked reference row. Rejoin its persistent guards. An uncommitted pair of provisional tails uses no reference positions. & At most \(D=d+4(q+20)\) old positions and two touched cycles, with at most \(v\) additional retreat to the next complete cycle. \\
Established contact & Restore the interval connecting the two fixed guards. Each branch pays for its own short detached end. & At most \(D\) old positions. The later free displacement calculation is unnecessary if the Grow repair fits; otherwise use the final-withdrawal row. \\
Lone attained overshoot & Restore the backed free prefix or the genuinely joined one. & The restored prefix reaches at least \(P+9v-D>P+8v\), before the next withdrawal. \\
Final withdrawal and normalization & Keep guards inside the target and switch to the in-home repair. Each old face was trimmed by ten sites. & At most two early Shrink inspections, included in the guard allowance. \\
Ordinary work & Use the old full target, its surviving guards, or the already observed service sample. & The sample has at least \(P-D>P/2\) matching sites. \\
\bottomrule
\end{longtable}

\subsection{Reader and reference estimates}

\begin{longtable}{@{}>{\raggedright\arraybackslash}p{.29\linewidth}>{\raggedright\arraybackslash}p{.63\linewidth}@{}}
\toprule
Quantity & Choice, after the data geometry has been fixed \\
\midrule
\endhead
Primitive pools & Thirteen wall probes, five mode probes, three groups of 59 replicas per reader. Thirty flags are needed to affect a decoded bundle. A reference group uses \(b_h+2\) probes and \(2v_h-1\) replicas. \\
Wall and mode rates & Wall refresh \(p^{12}\), new wall choice Bernoulli \(p\), CLEAN probability \(p^5\), RUN probability \(p^3(1-p^2)\). \\
Reader radius and buffer & \(W=\lceil p^{-3/2}\rceil\), \(\ell'_0=\lfloor p^{-1/2}\rfloor\), with the fixed extra boundary sites of Lemma~\ref{lem:reader}. \\
Reader time & \(B_0=\lceil p^{-14}\rceil\) sets the prefix, \(H=\lceil p^{-6}\rceil\) the suffix, and \(S=\lceil p^{-4}\rceil\) the time budget. A fully blanked segment gets at least \(S/2\) ticks to publish. \\
All-start mixing & Eight suffix slabs, enduring walls in the two exterior strips, and a CLEAN period of at least \(5W+10\) in each isolated segment. \\
Anchor and cutoff & \(A_n=8^{n+n_0}\), \(k_n=\lfloor A_n/8\rfloor\), \(R_n=2^{\lfloor A_n/4\rfloor}\), and the stripe width \(P^n\). \\
Short and long supports & Short radius \(p^{-16}\); long radius \(p^{-(b_h+5)}\), including the reference refresh prefix of length \(\lceil p^{-(b_h+3)}\rceil\). Leaf cutoff exponents \(17\) and \(b_h+6\) in the minimum side. \\
Tree and chain constants & Bounded child-pair count \(8(P+2)^2(T+2)^2\), enlarged-leaf overlap degree fixed by geometry, and the geometric count \(C_1 r_\lambda^4\) for each ordered ball. Their use precedes the final small-\(p\) choice. \\
\bottomrule
\end{longtable}

\section{Dependencies of the lemmas}\label{sec:dependencies}

The table is the exact directed \texttt{\string\ref} graph on the labelled lemmas and Proposition~\ref{prop:organization}, with self-references omitted, using each statement and its proof; the organization row includes its surrounding induction explanation and final assembly. It has 35 nodes and 90 edges. An arrow means that the row text references the result in its Uses column, including explanatory forward pointers. Used in is the exact reverse relation. References to definitions, equations, sections and theorems are outside this induced graph. Unlabelled native-rule and lower-rank premises remain as stated in each proof. The theorem-level arguments are also outside the graph: Section~\ref{sec:halting} uses Lemmas~\ref{lem:stripes}, \ref{lem:joint}, \ref{lem:locking} and~\ref{lem:ascent}, and Section~\ref{sec:gibbs} uses Lemmas~\ref{lem:past} and~\ref{lem:clock-witness} with Lemmas~\ref{lem:blocking} and~\ref{lem:doubling}.

At rank j, references to Proposition~\ref{prop:organization} within an induction step invoke a rank-\((j-1)\) clause or identify the current conclusion being proved; the assembly invokes the established rank-j lemmas. In the birth branch the computing test is proved before strong issue, incoming clearance and successor properness are used.

The supplied-interval Lemma~\ref{lem:local-wound} reproduces the first-clean-Purge count independently and uses only lower transitions and held-flank clearance. The held-flank proof fixed-backing pointer goes to the independent step~2 of Lemma~\ref{lem:ordinary-response}, so this longer textual cycle also supplies no missing premise.

Every pair of reciprocal textual references has the following interpretation.
\par\noindent Proposition~\ref{prop:organization} and Lemma~\ref{lem:ballot} are reciprocal only because the lemma invokes a lower-rank clause or names the current conclusion, while the assembly uses the completed lemma at rank j.
\par\noindent Proposition~\ref{prop:organization} and Lemma~\ref{lem:end-ledger} are reciprocal only because the lemma invokes a lower-rank clause or names the current conclusion, while the assembly uses the completed lemma at rank j.
\par\noindent The transition/ascent pair: the transition proof only points to the separately centered cap tests of Lemma~\ref{lem:ascent}; the ascent uses the transition conventions.
\par\noindent The flanks/ordinary-response pair reuses only the fixed-backing induction in step~2, which is closed independently of the local response or repair.
\par\noindent The flanks/seed pair concerns the paired-seam application conditional on an already supplied full seed; obtaining that seed uses only the independently proved backing and ordinary geometry.
\par\noindent The flanks/attachment pair points from the held-wall mechanism to its issued-source application; attachment supplies the walls before invoking clearance.
\par\noindent The ordinary-response/seed pair points forward to the count of open-end episodes; the ordinary proof establishes the nonlatching geometry before that independent counting application.
\par\noindent The ordinary-response/attachment pair is an application pointer from ordinary work; the local calculation is proved earlier as Lemma~\ref{lem:local-wound} without an issued-source conclusion.
\par\noindent The attachment/components pair uses only the independent clean-evolution paragraph in the forward direction; the component count receives its initial source interval from attachment.
\par\noindent The attachment/remote pair points to an independently rederived sterile-band induction; source attachment does not use the remote-source conclusion.
\par\noindent The competitors/calendar pair uses only the independent old middle-wall selection, which is made without births or the terminal-exclusion conclusion.
\par\noindent The calendar/end-ledger pair uses the \(97q\) guard only in the conditional filling remark; the end ledger selects that guard before invoking the clear calendar.

Longer textual cycles reduce to these pairs. Delete the pointer or independent-sub-argument direction of each pair (flanks to ordinary-response, seed and attachment; ordinary-response to seed and attachment; attachment to components and remote; competitors to calendar; calendar to end-ledger; transition to ascent) and every reference from a lemma of the one-rank induction to Proposition~\ref{prop:organization}, which invokes a rank-\((j-1)\) clause or names the conclusion currently being established. The remaining graph is acyclic. In particular the cycle attachment--remote--thin-density--attachment is broken at the attachment/remote pointer, the cycles through the rank-\((j-1)\) classification clause cited by Lemma~\ref{lem:ordinary-start} are broken at that clause, and every cycle through Lemma~\ref{lem:transition} is broken at its pointer to Lemma~\ref{lem:ascent}.

\begin{longtable}{@{}>{\raggedright\arraybackslash}p{.14\linewidth}>{\raggedright\arraybackslash}p{.40\linewidth}>{\raggedright\arraybackslash}p{.38\linewidth}@{}}
\toprule
Result & Uses & Used in \\
\midrule
\endhead
\ref{lem:clean} & --- & \ref{lem:past} \\
\ref{lem:sparse-witness} & --- & \ref{lem:joint} \\
\ref{lem:reader} & --- & \ref{lem:receiver-mixing}, \ref{lem:stripes} \\
\ref{lem:receiver-mixing} & \ref{lem:reader} & \ref{lem:stripes} \\
\ref{lem:stripes} & \ref{lem:reader}, \ref{lem:receiver-mixing} & \ref{lem:joint}, \ref{lem:past} \\
\ref{lem:joint} & \ref{lem:sparse-witness}, \ref{lem:stripes} & \ref{lem:clock-witness} \\
\ref{prop:organization} & \ref{lem:transition}, \ref{lem:ballot}, \ref{lem:classification}, \ref{lem:exclusion}, \ref{lem:end-ledger}, \ref{lem:advance} & \ref{lem:ordinary-start}, \ref{lem:ballot}, \ref{lem:attachment}, \ref{lem:end-ledger}, \ref{lem:ascent}, \ref{lem:clock-witness} \\
\ref{lem:transition} & \ref{lem:ascent} & \ref{prop:organization}, \ref{lem:ordinary-start}, \ref{lem:local-wound}, \ref{lem:startup}, \ref{lem:advance}, \ref{lem:two-period}, \ref{lem:ascent} \\
\ref{lem:ordinary-start} & \ref{prop:organization}, \ref{lem:transition} & \ref{lem:ordinary-response}, \ref{lem:seed}, \ref{lem:attachment}, \ref{lem:thin-density} \\
\ref{lem:flanks} & \ref{lem:ordinary-response}, \ref{lem:seed}, \ref{lem:attachment} & \ref{lem:local-wound}, \ref{lem:ordinary-response}, \ref{lem:seed}, \ref{lem:ballot}, \ref{lem:attachment}, \ref{lem:exclusion}, \ref{lem:end-ledger} \\
\ref{lem:local-wound} & \ref{lem:transition}, \ref{lem:flanks} & \ref{lem:ordinary-response}, \ref{lem:attachment} \\
\ref{lem:ordinary-response} & \ref{lem:ordinary-start}, \ref{lem:flanks}, \ref{lem:local-wound}, \ref{lem:seed}, \ref{lem:attachment} & \ref{lem:flanks}, \ref{lem:seed}, \ref{lem:ballot}, \ref{lem:attachment}, \ref{lem:thin-density} \\
\ref{lem:seed} & \ref{lem:ordinary-start}, \ref{lem:flanks}, \ref{lem:ordinary-response} & \ref{lem:flanks}, \ref{lem:ordinary-response}, \ref{lem:ballot}, \ref{lem:exclusion}, \ref{lem:end-ledger}, \ref{lem:two-period} \\
\ref{lem:ballot} & \ref{prop:organization}, \ref{lem:flanks}, \ref{lem:ordinary-response}, \ref{lem:seed} & \ref{prop:organization}, \ref{lem:attachment}, \ref{lem:inheritance}, \ref{lem:exclusion}, \ref{lem:cap-cut} \\
\ref{lem:classification} & --- & \ref{prop:organization}, \ref{lem:exclusion} \\
\ref{lem:attachment} & \ref{prop:organization}, \ref{lem:ordinary-start}, \ref{lem:flanks}, \ref{lem:local-wound}, \ref{lem:ordinary-response}, \ref{lem:ballot}, \ref{lem:components}, \ref{lem:remote} & \ref{lem:flanks}, \ref{lem:ordinary-response}, \ref{lem:components}, \ref{lem:thin-density}, \ref{lem:remote}, \ref{lem:calendar}, \ref{lem:end-ledger} \\
\ref{lem:components} & \ref{lem:attachment} & \ref{lem:attachment}, \ref{lem:thin-density} \\
\ref{lem:thin-density} & \ref{lem:ordinary-start}, \ref{lem:ordinary-response}, \ref{lem:attachment}, \ref{lem:components} & \ref{lem:remote}, \ref{lem:inheritance}, \ref{lem:competitors} \\
\ref{lem:remote} & \ref{lem:attachment}, \ref{lem:thin-density} & \ref{lem:attachment}, \ref{lem:inheritance}, \ref{lem:competitors} \\
\ref{lem:inheritance} & \ref{lem:ballot}, \ref{lem:thin-density}, \ref{lem:remote} & \ref{lem:exclusion}, \ref{lem:competitors}, \ref{lem:end-ledger} \\
\ref{lem:exclusion} & \ref{lem:flanks}, \ref{lem:seed}, \ref{lem:ballot}, \ref{lem:classification}, \ref{lem:inheritance} & \ref{prop:organization} \\
\ref{lem:endpoint} & --- & \ref{lem:end-ledger}, \ref{lem:advance}, \ref{lem:two-period} \\
\ref{lem:competitors} & \ref{lem:thin-density}, \ref{lem:remote}, \ref{lem:inheritance}, \ref{lem:calendar} & \ref{lem:calendar}, \ref{lem:end-ledger} \\
\ref{lem:calendar} & \ref{lem:attachment}, \ref{lem:competitors}, \ref{lem:end-ledger} & \ref{lem:competitors}, \ref{lem:end-ledger} \\
\ref{lem:end-ledger} & \ref{prop:organization}, \ref{lem:flanks}, \ref{lem:seed}, \ref{lem:attachment}, \ref{lem:inheritance}, \ref{lem:endpoint}, \ref{lem:competitors}, \ref{lem:calendar} & \ref{prop:organization}, \ref{lem:calendar} \\
\ref{lem:startup} & \ref{lem:transition} & \ref{lem:advance} \\
\ref{lem:advance} & \ref{lem:transition}, \ref{lem:endpoint}, \ref{lem:startup}, \ref{lem:two-period} & \ref{prop:organization} \\
\ref{lem:two-period} & \ref{lem:transition}, \ref{lem:seed}, \ref{lem:endpoint} & \ref{lem:advance}, \ref{lem:clock-witness} \\
\ref{lem:locking} & --- & \ref{lem:cap-cut}, \ref{lem:past}, \ref{lem:clock-witness} \\
\ref{lem:cap-cut} & \ref{lem:ballot}, \ref{lem:locking} & \ref{lem:ascent} \\
\ref{lem:ascent} & \ref{prop:organization}, \ref{lem:transition}, \ref{lem:cap-cut} & \ref{lem:transition} \\
\ref{lem:past} & \ref{lem:clean}, \ref{lem:stripes}, \ref{lem:locking} & \ref{lem:clock-witness} \\
\ref{lem:clock-witness} & \ref{lem:joint}, \ref{prop:organization}, \ref{lem:two-period}, \ref{lem:locking}, \ref{lem:past} & --- \\
\ref{lem:blocking} & --- & --- \\
\ref{lem:doubling} & --- & --- \\
\bottomrule
\end{longtable}

